\section{Conclusiones}
\label{sec:conclusiones}

Este trabajo presentó el diseño, implementación y evaluación empírica del Sistema de Gestión de
Pre-Sustentaciones UTEQ, siguiendo Design Science Research \cite{peffers2007} de forma instanciada y
verificable. Los 12 requisitos funcionales están implementados y el flujo académico completo se verificó
funcionando end-to-end contra una topología de producción real. La evaluación empírica combina
estadística descriptiva e inferencial genuina --- el hallazgo más sólido es la reducción de latencia de
13.9$\times$ por efecto de caché Redis, con significancia estadística fuerte ($p < 10^{-10}$) y tamaño
de efecto máximo --- y auditoría de seguridad multi-herramienta que confirma ausencia de inyección SQL en
el código propio; se corrigió la única vulnerabilidad de dependencias que afectaba código servido al
navegador (\texttt{@angular/core}, detectada por ZAP), y quedan abiertas 14 vulnerabilidades en
herramientas de build que no se envían a producción.

La contribución más distintiva de este informe no es técnica sino metodológica: el proceso de
elaboración de este mismo proyecto incluyó la detección real de datos académicos fabricados en fases
anteriores (usabilidad, rendimiento web, evidencia de trazabilidad, integridad de configuración de
despliegue, custodia de consentimientos informados) y su corrección sistemática, documentada de forma
trazable en \texttt{docs/mediciones/DATA-PROVENANCE.md} y en la bitácora de cada fase del proyecto. Esa
corrección no fue cosmética: implicó reportar una cobertura de código de 38.88\,\% en vez de un
``$>$60\,\%'' previamente afirmado sin evidencia, y declarar la usabilidad como ``sin datos'' en vez de
mantener un 91.25/100 fabricado.

Las brechas que quedan (cobertura de pruebas incompleta, ausencia de datos SUS reales, despliegue público
pendiente, dependencias vulnerables sin corregir) se declaran explícitamente en la
Sección~\ref{sec:trabajo-futuro} en vez de ocultarse o rellenarse con contenido que simule completitud.
Esta decisión --- priorizar un informe más corto pero verificable sobre uno más extenso pero con
contenido fabricado --- es, en sí misma, la conclusión metodológica central de este trabajo: en un
proyecto con antecedentes reales de fabricación de evidencia, la integridad de lo que se reporta importa
más que la completitud aparente de lo que se reporta.
\newpage
